%%%%%%%%%%  scan idx 321  |  folio 308  %%%%%%%%%%
sa valeur sur l'élément \(1+Xt\) de \(\Ghat\circ i(K[X])\), est injective; son image est le sous-groupe de \(H(K[X])\) formé des éléments \(\theta\) tels que (avec des notations analogues à celles de 1.4 et 1.5)~:
\[
\begin{array}{ll}
\text{i)}& H(f_1)(\theta)=0;\\[0.3em]
\text{ii)}& \forall n\geq1,\quad
H(f_n)\bigl(u_H^n(\theta)\bigr)\equiv u_H^n(\theta)
\ \mathrm{mod.}\ F^nH(K[X_1,\ldots,X_n]).
\end{array}
\]

Le théorème résulte immédiatement du théorème 1.6 et du lemme 1.9, compte tenu de ce que \(H\circ g(K[X])\) ne dépend pas de la filtration choisie sur \(K[X]\).

Le théorème analogue obtenu en remplaçant \(\Ghat\) par \((\Ghat)^k\) est également vrai (cf. 1.7)

Quand il n'y aura pas de risque de confusion, nous noterons le foncteur \(\Ghat\circ i\) par \(\Ghat\), ce qui est légitime, car \(i\) identifie \(\mathcal C_0\) à une sous-catégorie de \(\mathcal C\).

\textbf{1.11.} On notera \(\mathcal G\) la catégorie des \(K\)-algèbres graduées, et \(j\) le foncteur de \(\mathcal G\) dans \(\mathcal C\) qui à une \(K\)-algèbre graduée \(B\) associe la \(K\)-algèbre \(B\) munie de la filtration canoniquement définie par sa graduation~; le foncteur «~gradué associé~»~: \(\mathcal C\to\mathcal G\), sera noté \(\operatorname{Gr}\). Il existe un isomorphisme de foncteurs \(\varepsilon:\operatorname{Id}_{\mathcal G}\xrightarrow{\sim}\operatorname{Gr}\circ j\).

Soit \(B\in\operatorname{Ob}(\mathcal G)\)~; on notera \(1+\widehat B^+\) le groupe multiplicatif des éléments de \(\prod\limits_{i=0}^{\infty}B^i\) d'augmentation \(1\). En le munissant de la filtration analogue à celle de \(\Ghat\), on définit ainsi un foncteur~: \(\mathcal G\to\Abfil\).

Par ailleurs, on définit pour tout \(A\in\operatorname{Ob}(\mathcal C)\) un homomorphisme de groupes filtrés~: \(\Ghat(A)\longrightarrow 1+\widehat{\operatorname{Gr}A}^{+}\), en associant à la série \(1+\ldots+a_nt^n+\ldots\), 
